% Pista alg-25i-q4b · Àlgebra
% Procedència: PAU Catalunya 2025 · Convocatòria de juny (incidències) · Sèrie 4 · Exercici 4 — Opció B
\documentclass[11pt,a4paper]{article}
\usepackage{pista}
\pistainfo{Catalunya 2025}{Convocatòria de juny (incidències)}{Sèrie 4}

\begin{document}

\section*{\textcolor{blauPista}{Exercici 4 --- Opció B}}

\noindent Considereu les matrius
\[
A = \begin{pmatrix} a & 0 \\ 1 & a^{2} \end{pmatrix}, \text{ on } a \text{ és un paràmetre real, i } B = \begin{pmatrix} 2 & 0 \\ 1 & 4 \end{pmatrix}.
\]

\subsection*{\textit{a)} \normalfont Trobeu els valors de $a$ que fan que les matrius $A$ i $B$ commutin. (Dues matrius $A$ i $B$ commuten si $AB = BA$.) \hspace*{\fill}\textnormal{[0,75~punts]}}

\pista{En primer lloc, calcula els dos productes, $A \cdot B$ i $B \cdot A$. Obtindràs dues matrius que depenen del paràmetre $a$.

\medskip
A continuació, iguala la matriu $A \cdot B$ amb la matriu $B \cdot A$, element per element, i obtindràs quatre equacions. No és tan difícil com sembla, perquè algunes d'aquestes equacions es compleixen sempre. Finalment, resol les que quedin.}

\subsection*{\textit{b)} \normalfont Trobeu els valors de $a$ per als quals la matriu $A$ és invertible. \hspace*{\fill}\textnormal{[0,75~punts]}}

\pista{Calcula el determinant de la matriu $A$, que dependrà del paràmetre $a$. Si el determinant és diferent de zero, significa que la matriu $A$ és invertible.

\medskip
Per tant, busca per a quins valors d'$a$ el determinant s'anul·la. La teva resposta són, precisament, tots els altres valors. Per exemple, si el valor que anul·la el determinant fos $a = 3$, llavors tu hauries d'escriure que la resposta és $a \neq 3$.}

\subsection*{\textit{c)} \normalfont Trobeu els valors de $a$ per als quals $A^{-1} = A$. \hspace*{\fill}\textnormal{[1~punt]}}

\pista{En lloc de calcular la inversa, multiplica els dos costats de la igualtat $A^{-1} = A$ per la matriu $A$. Com que $A \cdot A^{-1} = I$, obtindràs una igualtat sense inverses:
\[
A \cdot A = I
\]
A continuació, calcula $A \cdot A$ i iguala-la amb la matriu identitat $I$, element per element.

\medskip
\textbf{Atenció:} obtindràs diverses equacions, i és important que el valor d'$a$ les compleixi totes alhora.}

\end{document}
